\begin{figure}[H]
  \centering
  \begin{tikzpicture}
    \begin{axis}[paperaxis,
      xlabel={Année}, ylabel={Indice, 2017 $= 100$},
      xtick={2017,2018,2019,2020,2021,2022},
      xticklabel style={/pgf/number format/1000 sep=},
      legend pos=north west, legend cell align=left,
    ]
      % Dividendes des menages : la serie focale, en bleu France.
      \addplot[franceblue, line width=1.5pt, mark=*, mark options={scale=0.8}]
        table[col sep=comma, x=year, y=dividendes] {figures/decouplage_fr.csv};
      \addlegendentry{Dividendes des ménages}
      % Profit et activite : les references, en gris ; trait plein le profit, tirete l'activite.
      \addplot[controlgrey, line width=1.1pt, mark=square*, mark options={scale=0.6}]
        table[col sep=comma, x=year, y=profit] {figures/decouplage_fr.csv};
      \addlegendentry{Profit opérationnel (EBE des SNF)}
      \addplot[controlgrey, line width=1.1pt, densely dashed, mark=triangle*, mark options={scale=0.7}]
        table[col sep=comma, x=year, y=pib] {figures/decouplage_fr.csv};
      \addlegendentry{Activité (PIB en volume)}
      % Reforme de 2018, en vert.
      \draw[baisse2018, dashed, line width=1pt]
        ({axis cs:2018,0}|-{rel axis cs:0,0}) -- ({axis cs:2018,0}|-{rel axis cs:0,1});
    \end{axis}
  \end{tikzpicture}
  \caption{Dividendes, profit et activité en France, indice base 100 en 2017}
  \label{fig:decouplage}
  \fignote{Trois séries françaises indexées à 100 en 2017 : les dividendes reçus par les ménages,
  poste D42 des comptes de secteur ; l'excédent brut d'exploitation des sociétés non financières, le
  profit opérationnel ; et le produit intérieur brut en volume, l'activité. En bleu la série focale,
  les dividendes ; en gris les deux grandeurs de référence, trait plein le profit, trait tireté
  l'activité. La ligne verte verticale marque la réforme de 2018. Les dividendes retombent en 2020
  puis bondissent en 2021 à 187, très au-dessus du profit, à 119, et de l'activité, à 103, qui ne
  font que revenir vers leur niveau de 2019. L'écart aux pays comparables en 2021 va dans le même
  sens : dividendes $+80$ points d'indice, profit $+6$, activité nul. Deux réserves de périmètre : les
  dividendes des ménages incluent les versements reçus de l'étranger et des fonds, pas seulement des
  sociétés non financières françaises, et l'excédent brut d'exploitation est un profit opérationnel
  avant intérêts et impôts, non le résultat net distribuable. La figure montre donc que les dividendes
  reçus par les ménages se découplent de l'activité et du profit opérationnel domestiques, non que les
  entreprises distribuent sans profit. Source : Eurostat.}
\end{figure}
